\chapter{Úvod}
\label{chap:introduction}

Pochopení datových struktur a algoritmů představuje důležitý základ pro každého studenta informatiky.
Tyto koncepty formují analytické myšlení budoucích softwarových inženýrů a jsou nezbytné pro návrh efektivních a škálovatelných programů.
Přesto se jedná o oblast, která studentům tradičně působí obtíže; algoritmy jsou ze své podstaty dynamické procesy plné abstraktních stavů a manipulací s pamětí, které je obtížné plně pochopit pouze ze statických textů, skript či kreseb na tabuli.

Ačkoliv v dnešní době existuje na internetu řada vizualizačních nástrojů, studenti i vyučující často naráží na problémy: roztříštěnost a nekompatibilitu.
Mnoho existujících aplikací je zastaralých, chybí jim responzivita, nepodporují český jazyk, nebo používají definice a pseudokódy, které se drobně či dokonce zásadně liší od materiálů probíraných na konkrétních univerzitách.
Tato nejednotnost má za následek stav, kdy studenti musí neustále přepínat mezi různými zdroji a porovnávat jejich informace se skripty, což je spíše matoucí než prospěšné pro jejich učení.
Chybí zde moderní, sjednocující výuková aplikace, která by byla nejen vizuálně atraktivní, ale především by mohla být oficiálně integrována do výuky a sladěna s konkrétními univerzitními skripty.

Hlavním cílem této diplomové práce je proto návrh a implementace interaktivní výukové webové aplikace, která tento problém řeší.
Aplikace si klade za cíl poskytnout studentům vizuální a intuitivní prostředí pro zkoumání datových struktur a třídicích algoritmů.
Nabízí unikátní kombinaci pohledů -- interaktivní manipulaci se stromy, obecné vizualizace řazení prvků nebo detailní krokování rekurzivních algoritmů s vazbou na dynamický pseudokód, což je kombinace funkcionalit, která u mnoha existujících řešení chybí.

Aplikace je primárně cílena na studenty univerzitních oborů informatiky jakožto podpůrný nástroj k přednáškám a cvičením.
Její moderní uživatelské rozhraní a přístupnost jsou navrženy tak, aby bariéra vstupu byla co nejnižší.
Sloužit tak může i středoškolským studentům, samoukům nebo komukoliv, kdo má zájem rozšířit své porozumění fungování základních počítačových algoritmů.

Text práce je logicky členěn do několika kapitol.
Po úvodu následuje \customref{Kapitola}{chap:content}, která shrnuje teoretické základy implementovaných datových struktur a algoritmů.
\customref{Kapitola}{chap:other_apps} nabízí rešerši a kritickou analýzu existujících vizualizačních řešení na trhu.
Na základě zjištěných nedostatků jsou v \customref{Kapitole}{chap:requirements} specifikovány požadavky na novou aplikaci.
\customref{Kapitola}{chap:technologies} se věnuje výběru vhodných moderních webových technologií a strategií renderování.
\customref{Kapitola}{chap:design} detailně popisuje softwarový návrh, informační architekturu a uživatelské rozhraní, zatímco \customref{Kapitola}{chap:implementation} rozebírá samotnou implementaci, architektonická řešení a překonávání technologických limitů.
Práce je zakončena \customref{Kapitolou}{chap:conclusion}, která shrnuje dosažené výsledky a nastiňuje možnosti budoucího rozvoje aplikace.
